\documentclass[11pt]{amsart}

\usepackage{amsmath}
\usepackage{amssymb}

%% ------------------------------------------------------------------
%% Theorem environments
%% ------------------------------------------------------------------
\theoremstyle{plain}
\newtheorem{theorem}{Theorem}
\newtheorem{lemma}{Lemma}[section]
\newtheorem{proposition}{Proposition}[section]

% Named results, cited by letter throughout, as in the manuscript.
\newtheorem*{theoremmain}{Theorem 1}
\newtheorem*{lemmaB}{Lemma B}
\newtheorem*{lemmaU}{Lemma U (rank-1 reduction)}
\newtheorem*{lemmaAprime}{Lemma A$'$ (generators along the cycle)}
\newtheorem*{propositionF}{Proposition F (forward window)}
\newtheorem*{lemmaC}{Lemma C (backward window)}

\theoremstyle{remark}
\newtheorem{remark}{Remark}[section]

%% ------------------------------------------------------------------

\begin{document}

\title{Pure cubic Thue equations are solvable in deterministic
  exponential time}

\author{The diophantine-fable expedition}
\thanks{Draft --- authorship placeholder.}

\date{Draft of 2026-07-22 (v2); all corrections from the independent
  verification passes applied; the numeric evaluation of one absolute
  constant is deferred and flagged (Remark~\ref{rem:cchain}).}

\begin{abstract}
We prove that pure cubic Thue equations $x^3 - d\,y^3 = m$ ($d \ge 2$
cube-free, not a cube; $m \ne 0$; input size $s = O(\log d + \log|m|)$
in binary) are decided, with full solution lists, by a deterministic,
unconditional algorithm in time $2^{O(s)}$. No worst-case bound of this
shape appears in the literature: the best effective height bounds
(Bugeaud--Gy\H{o}ry) are exponential in the coefficient
\emph{bit-size}, so enumeration is doubly exponential, and the closest
prior (Smart, ANTS-II 1996) is an explicitly practical per-method
estimate, not a worst-case bound. The method walks the rank-one unit orbit of
norm-$m$ representatives from the Buchmann--Williams infrastructure and
confines Thue-shape indices to a two-sided window: an elementary $O(1)$
forward bound, and a backward bound quasi-linear in the regulator $R$,
via Sha's Matveev-based endgame re-run with the true root-ratio gap
$3R/2$, against which the $R$ in Matveev's height parameter cancels. A
prototype matches PARI's certified \texttt{thue} on 1120/1120 grid
instances; an independent reimplementation matches 13/13, including
three high-index fields. All constants are explicit modulo
one flagged final evaluation.
\end{abstract}

\maketitle

%% ==================================================================
\section{Introduction}\label{sec:intro}
%% ==================================================================

\subsection{Smale's fifth problem and the Thue gap}\label{sec:smale}

Smale's fifth problem (\emph{Mathematical problems for the next
century} \cite{Sma98}; see also Lagarias's quotation of the AMS reprint
in \cite{Lag79}, \S 1.4) asks for an algorithm deciding, given
$f \in \mathbb{Z}[u,v]$, whether $f(x,y) = 0$ has a solution in
$\mathbb{Z}^2$, in time $2^{s^c}$ for a universal constant $c$ ---
exponential time $2^{\mathrm{poly}(s)}$ in the input size $s$. The size
measure is \emph{dense}:
$s(f) = \sum_{|\alpha| \le \deg f} \max(\log|a_\alpha|, 1)$, so every
coefficient slot up to the total degree contributes at least $1$. Smale
notes that even decidability of the general problem is open.

Within this landscape, Thue equations $F(x,y) = m$ with $F$ an
irreducible binary form of degree $\ge 3$ occupy a deceptively
classical-looking stratum. Their finiteness is Thue's theorem; their
effectivity is Baker's. Yet the best general effective height bounds
--- Bugeaud--Gy\H{o}ry \cite{BG96} --- bound $\log\max(|x|,|y|)$
\emph{polynomially in the coefficient magnitude} $H$, that is,
\textbf{exponentially in the coefficient bit-size} $\log H$.
Enumeration up to such a bound therefore costs doubly exponential time
in $s$, and the following gap results:

\begin{quote}
Even Thue equations of \textbf{fixed degree 3} are not known to be
decidable in time $2^{s^c}$ by the height-bound route.
\end{quote}

This paper closes that gap for the pure cubic family. Throughout,
$d \ge 2$ is cube-free and not a perfect cube, $m \ne 0$, and
$s = O(\log d + \log |m|)$ denotes the bit-size of the input $(d, m)$
in binary.

\subsection{Prior work: effectivity and practicality, but no worst-case
  bound}\label{sec:prior}

The algorithmic Thue literature is rich but states no worst-case
complexity theorem of the shape we prove.

\begin{itemize}
\item \textbf{Tzanakis--de Weger} \cite{TdW89} set up the now-standard
  practical pipeline (Baker bound, LLL-based reduction,
  small-solutions search). Their paper states no complexity bound; the
  necessary algebraic data --- fundamental units and the factorization
  of $m$ in the field --- are \emph{explicitly assumed known}.
\item \textbf{Bilu--Hanrot} \cite{BH96} extend the pipeline to high
  degree (replacing the $r$-dimensional LLL reduction by
  continued-fraction reduction of two-term forms). Again: no
  complexity bound; the phrase is ``solve in reasonable time \dots\
  provided the necessary algebraic number theory data is available''.
\item \textbf{Smart} \cite{Sma96} is the closest prior and the work we
  position against. It is a complexity analysis of the Tzanakis--de
  Weger method at fixed degree $n$, measured in $\log|m| + \log L(F)$
  ($L(F)$ the sum of the coefficient moduli), and it is explicitly an
  analysis of practical difficulty --- ``I know of no analysis as to
  how difficult it is in practice to actually find all the
  solutions'' \cite[p.~363]{Sma96}. The solution space is divided into
  small solutions (direct search), medium solutions
  (continued-fraction reduction), and large solutions (eliminated by
  linear forms in logarithms), and the announced outcome is ``an
  overall exponential complexity bound in terms of $\log L(F)$''
  \cite[p.~364]{Sma96} --- exponential in the input size --- with the
  practical failure modes stated explicitly: large $m$, large
  regulator or class number, near-coincident roots of
  $F(x,1)$.\footnote{Verified against the openly available portion of
    the primary text: the publisher's two-page preview of \cite{Sma96}
    (pp.~363--364) confirms the fixed-degree setting, the complexity
    measure $\log|m| + \log L(F)$, the three-phase division of the
    solution space, the failure modes, and the announced ``overall
    exponential complexity bound in terms of $\log L(F)$''; the
    concluding quotation is reproduced verbatim in R.~J. Stroeker's
    zbMATH review (Zbl 0896.11009). The remainder of the text
    (pp.~365--373) is paywalled; per-phase exponent shapes reported in
    secondary sources --- e.g.\ a small-solutions cost of
    $O(|m|^{1/(n-2)})$ --- are therefore not quoted above, and nothing
    in this section depends on them: the gap claim rests on the sweep
    of the books and surveys, not on the exact exponents of
    \cite{Sma96}.} Smart's conclusion, quoted verbatim in Stroeker's
  review of the paper (Zbl 0896.11009): ``So although solving Thue
  equations is now considered trivial we seem to be a long way off
  from a `good' algorithm''. It is a per-method practical estimate ---
  not a rigorous worst-case decision theorem, and no bound of the
  shape of Theorem~\ref{thm:main} appears there.
\item \textbf{PARI/GP} \cite{PARI} prints one genuine unconditional
  complexity statement for \texttt{thue}: when $F(x,1)$ has \emph{no
  real root} the equation can be solved unconditionally in time
  $|m|^{1/\deg F}$. Every pure cubic has a real root, so even this
  rank-limited printed bound does not cover our family; and
  \texttt{thue}'s output is GRH-conditional unless certified.
\end{itemize}

A systematic literature verification across the books and surveys
(Smart's 1998 book, Hanrot's LLL survey, Evertse--Gy\H{o}ry,
Waldschmidt, Ga\'al, Kim, Gherga--Siksek) found effectivity or
practicality statements only.
The theorem below therefore appears to be a citable gap in the
literature, positioned --- as we have just done --- against
\cite{Sma96}.

\subsection{The main theorem}\label{sec:mainthm}

\begin{theorem}[Main Theorem]\label{thm:main}
There is a deterministic, unconditional algorithm that, given cube-free
$d \ge 2$ (not a cube) and $m \ne 0$ in binary, with input size
$s = O(\log d + \log|m|)$, decides whether
\[
x^3 - d\,y^3 = m
\]
has a solution in $\mathbb{Z}^2$ --- and lists all solutions --- in
time $2^{O(s)}$.
\end{theorem}

The algorithm is fully specified and its correctness proof is complete;
every constant in the analysis is effective and explicitly computable.
The single item deferred --- flagged wherever it occurs --- is the
\emph{numeric evaluation} of one absolute constant chain
(Remark~\ref{rem:cchain}): the theorem holds as stated, with the window
bound of Lemma C semi-explicit, i.e.\ explicit modulo final constant
evaluation.

\subsection{Proof strategy: seven steps}\label{sec:strategy}

Let $K = \mathbb{Q}(\alpha)$, $\alpha = \sqrt[3]{d}$, a complex cubic
field with unit group $\langle -1, \varepsilon\rangle$ of rank one. A
solution $(x, y)$ is an element $\gamma = x - y\alpha$ of norm $m$; all
norm-$m$ elements fall into finitely many orbits under the unit action;
and ``being of Thue shape'' is the vanishing of one coordinate --- the
$\alpha^2$-coordinate --- along the orbit, which is an integer linear
recurrence of order 3 in the orbit index $k$. The algorithm
(Section~\ref{sec:algorithm}):

\begin{enumerate}
\item compute the regulator, the Voronoi cycle, and the fundamental
  unit by the Buchmann--Williams rank-one infrastructure ---
  deterministic, unconditional, no class group needed;
\item enumerate the ideals of norm $|m|$ in $O_K$ and extract a
  generator for each principal one, by the same cycle walk
  (Lemma A$'$);
\item unit-reduce each generator (Lemma U);
\item set up the order-3 integer recurrences (forward and reversed)
  for the cleared $\alpha^2$-coordinate;
\item bound the forward window elementarily: vanishing indices
  $k \le 3$ (Proposition F);
\item bound the backward window by one application of Matveev's
  theorem with the true root-ratio gap $3R/2$: vanishing indices
  $-k \le N_C = \widetilde{O}(R + s)$ (Lemma C);
\item walk the windows with exact integer arithmetic, collect the
  zeros, recover $(x,y)$, and apply the mandatory final filter
  $x^3 - dy^3 = m$.
\end{enumerate}

The technical heart is Step 6, and it turns on the guiding principle
of this program, stated here once and referred forward to
Remark~\ref{rem:blackbox}: \textbf{always ask exponential in what} ---
every bound must be calibrated against the \emph{input} size $s$,
never against an auxiliary quantity such as the regulator, for the
regulator itself is exponential in the input, $R = 2^{\Theta(s)}$
(Section~\ref{sec:regulator}). A black-box application of Sha's
explicit Skolem-type theorems \cite{Sha19} gives a window of size
$e^{22.5R}$ --- which is $2^{O(s)}$ only under a \emph{redefined} size
measure containing $R$ itself; calibrated against the input it is
doubly exponential (Remark~\ref{rem:blackbox} preserves this near-miss
as a cautionary lesson). The repair is not a
new theorem of transcendence theory but a re-run of Sha's own endgame
with the instance's true data: the two maximal-modulus roots of the
reversed recurrence are the conjugate complex pair, the third root is
smaller by the exact factor $e^{3R/2}$, and in the resulting
inequality the $R$ inside Matveev's height parameter $A_3 = 4R$
\textbf{cancels} against the gap $3R/2$ --- leaving a window
quasi-linear in $R$, hence $2^{O(s)}$ in honest input size.

\subsection{Deciding without exhibiting: the Pell analogy}
\label{sec:pell}

The philosophy is the one the Pell equation teaches one rung down. The
fundamental solution of $x^2 - Dy^2 = 1$ has $\Theta(\sqrt{D})$ digits
--- exponential in the bit-size of $D$ --- yet solvability questions
about generalized Pell equations are decided by walking orbits and
comparing positions, never by expanding the fundamental solution
unless it is actually needed. One rung up, the obstruction is steeper:
by dense size $s = 20$ the fundamental unit of
$\mathbb{Q}(\sqrt[3]{d})$ already reaches \textbf{1593 digits} per
coefficient (Section~\ref{sec:unitgrowth}, all records certified),
against roughly 145 digits for the Pell envelope at the same size. The
infrastructure machinery of Buchmann--Williams \cite{BW88} keeps every
\emph{per-step} object at polynomial size --- minima are represented
by poly-size ideal data, ``whereas the order of magnitude of a minimum
can be as large as $\exp\sqrt{D}$'' \cite[Prop.~2.11]{BW88} --- and
full expansion is invoked only where the $2^{O(s)}$ budget explicitly
covers it. Stated precisely, the philosophy is that the algorithm
never \emph{enumerates solutions} --- objects whose digit counts are
exponential in $s$ --- not that it never writes down a large number.
At the exponential-time target of this paper, \emph{expanding} the
fundamental unit and the Binet data is a legitimate simplification,
adopted deliberately in Lemma A$'$: plain exact products along the
cycle suffice. Compact representations become essential only for the
stronger targets --- poly-size certificates, NP membership ---
discussed in Remark~\ref{rem:beyond}.

\subsection{Organization}\label{sec:organization}

Section~\ref{sec:prelim} fixes notation and collects the
field-theoretic preliminaries, including the two regulator floors the
proof leans on. Section~\ref{sec:algorithm} states the algorithm.
Section~\ref{sec:lemmas} proves the four supporting results (Lemmas B,
U, A$'$, C and Proposition F), with all corrections from independent
verification applied. Section~\ref{sec:proofmain} assembles the proof of
Theorem~\ref{thm:main}. Section~\ref{sec:companion} reports the
computational companions. Section~\ref{sec:remarks} collects remarks
and open problems.

%% ==================================================================
\section{Preliminaries}\label{sec:prelim}
%% ==================================================================

\subsection{The field, its embeddings, and its units}\label{sec:field}

Fix cube-free $d \ge 2$, not a cube, and write $d = ab^2$ with $a, b$
coprime and squarefree. Let $\alpha = \sqrt[3]{d}$ (real) and
$K = \mathbb{Q}(\alpha)$, a cubic field of signature $(1,1)$: one real
embedding $\sigma_1$ and one conjugate pair $\sigma_2, \sigma_3 =
\bar\sigma_2$, with $\sigma_j(\alpha) = \omega^{j-1}\alpha$ for
$\omega = e^{2\pi i/3}$. The discriminant is $-27a^2b^2$ if
$d \not\equiv \pm 1 \pmod 9$ and $-3a^2b^2$ if
$d \equiv \pm 1 \pmod 9$; in particular, for every pure cubic field
\[
|\mathrm{disc}(K)| \ge 108,
\]
the minimum being attained at $d = 2$. We will use this floor twice.

The unit group is $O_K^\times = \langle -1, \varepsilon\rangle$ of
rank one. The torsion subgroup is $\{\pm 1\}$: any root of unity in
$K$ is fixed by the real embedding $\sigma_1$, hence equals $\pm 1$.
Normalize the fundamental unit by
\[
\lambda := \sigma_1(\varepsilon) > 1, \qquad R := \log\lambda,
\]
which is always possible (replace $\varepsilon$ by
$\pm\varepsilon^{\pm1}$). Write $\mu = \sigma_2(\varepsilon)$; then
$|\mu| = |\bar\mu| = \lambda^{-1/2}$, so the embedding moduli of
$\varepsilon$ are $(\lambda, \lambda^{-1/2}, \lambda^{-1/2})$ and $R$
is the regulator of $K$.

\begin{lemma}[$N(\varepsilon) = +1$ always]\label{lem:normone}
With the normalization $\lambda > 1$, the fundamental unit of a
complex cubic field has norm $+1$.
\end{lemma}

\begin{proof}
$N(\varepsilon) = \sigma_1(\varepsilon)\,|\sigma_2(\varepsilon)|^2
= \lambda|\mu|^2 > 0$, and a unit norm is $\pm 1$.
\end{proof}

This one-liner is load-bearing: it makes the \emph{reversed} orbit
recurrence integral (Step 4).

\subsection{The index formula and the clearing constant $3b$}
\label{sec:index}

By Dedekind \cite{Ded00} (see \cite{AW04} for a standard modern
account), for $d = ab^2$ cube-free,
\[
[\,O_K : \mathbb{Z}[\alpha]\,] \;=\; b \cdot
\begin{cases}
3 & \text{if } d \equiv \pm 1 \pmod 9,\\
1 & \text{otherwise.}
\end{cases}
\]
Consequently the index always divides $3b$, and for every $g \in O_K$,
writing $g = c_0(g) + c_1(g)\alpha + c_2(g)\alpha^2$ with
$c_i(g) \in \mathbb{Q}$, the cleared coordinates $3b\,c_i(g)$ are
rational integers. All integer recurrences below are stated for the
$3b$-cleared $\alpha^2$-coordinate. We emphasize --- because a draft
of this work got it wrong, and empirical testing reveals that the
naive ``clear by 3'' recipe fails at 91 values of $d < 500$ (e.g.\
$d = 12, 45, 175$) --- that clearing by $3$ alone is \textbf{not}
sufficient: the correct clearing constant is $3b$.

\subsection{Regulator bounds, above and below}\label{sec:regulator}

\textbf{Above.} By Lenstra \cite[Thm.~6.5]{Len92}, $hR \le
|\Delta|^{1/2}\cdot \mathrm{polylog}|\Delta|$ at fixed degree; since
$|\Delta| \le 27d^2 = 2^{O(s)}$, we get
\[
R \;\le\; hR \;=\; 2^{O(s)}.
\]
This is the honest calibration to keep in view throughout: the
regulator is \emph{exponential} in the input size, so any window or
cost ``polynomial in $R$'' is exponential in $s$, and anything
``exponential in $R$'' is doubly exponential in $s$.

\textbf{Below, twice.} Astudillo--D\'iaz y D\'iaz--Friedman
\cite[Thm.~7 and Table~1]{ADF16} prove the sharp lower bound for
complex cubic fields:
\[
R \;\ge\; 0.2811995743\ldots,
\]
attained uniquely at the field of discriminant $-23$ (defined by
$x^3 - x - 1$; the fundamental unit is the plastic number). Moreover,
with exactly three exceptions --- the fields of discriminant $-23,
-31, -44$, with regulators $0.2811\ldots$, $0.3822\ldots$,
$0.6093\ldots$ --- every complex cubic field satisfies $R > 0.79$.
Since pure cubic fields have $|\mathrm{disc}| \ge 108$, \textbf{none
of the three exceptional fields is pure}, so:
\[
\text{every pure cubic field satisfies } R > 0.79.
\]
This floor is \textbf{load-bearing}: Matveev's theorem imposes the
side condition $A_3 \ge |\log(\bar\mu/\mu)|$, and
$|\log(\bar\mu/\mu)|$ can approach $\pi$. Our choice $A_3 = 4R$
satisfies the side condition precisely because $4R \ge 3.16 > \pi$ ---
a margin of $0.018$. At the overall complex-cubic minimum $R = 0.2812$
the condition would \emph{fail} ($4R = 1.12 < \pi$); the theorem
survives on the pure-cubic discriminant floor by less than two
hundredths. (The forward window, by contrast, needs only the
unrestricted floor $R \ge 0.28119$.)

\subsection{Heights}\label{sec:heights}

$h$ denotes the absolute logarithmic Weil height. We use throughout:

\begin{itemize}
\item $h$ is Galois-invariant: $h(\sigma_j(\gamma)) = h(\gamma)$.
\item $|\log|\beta|| \le [\mathbb{Q}(\beta):\mathbb{Q}]\cdot h(\beta)$
  for $\beta \ne 0$.
\item $h(\varepsilon) = R/3$: the Mahler measure of the minimal
  polynomial of $\varepsilon$ is $\lambda$ (the complex place, of
  local degree 2, contributes nothing since $|\mu| < 1$), so
  $h(\varepsilon) = \tfrac13\log\lambda$. Note the complex place
  carries weight 2 in all such computations; in particular
  $h(\bar\mu/\mu) \le 2h(\varepsilon) = 2R/3$.
\item $h(\omega) = 0$, $h(3d^{2/3}) = \log 3 + \tfrac23\log d$.
\end{itemize}

\subsection{The coordinate extraction and the orbit recurrence}
\label{sec:coordinate}

For $g \in K$ the discrete-Fourier identity of the power basis reads
\[
\sigma_1(g) + \omega\,\sigma_2(g) + \omega^2\sigma_3(g) \;=\;
3d^{2/3}\,c_2(g),
\]
so the $\alpha^2$-coordinate is a linear combination of the embeddings
with weights of modulus $1/(3d^{2/3})$. An element $\gamma \ne 0$ and
the orbit $g_k = \gamma\varepsilon^k$ give the Binet form
\[
c_2(\gamma\varepsilon^k) \;=\; b_1\lambda^k + b_2\mu^k
+ \bar b_2\bar\mu^k,
\qquad b_1 = \frac{\sigma_1(\gamma)}{3d^{2/3}},\quad
b_2 = \frac{\omega\,\sigma_2(\gamma)}{3d^{2/3}},
\]
and \textbf{all Binet coefficients are nonzero} whenever
$\gamma \ne 0$ (the embeddings of a nonzero element are nonzero, and
the DFT weights have modulus $1/(3d^{2/3}) \ne 0$). This verifies,
once and for all, the implicit ``all $b_j \ne 0$'' hypothesis under
which the Skolem-type machinery of Section~\ref{sec:lemmas} operates.

Since $[K:\mathbb{Q}] = 3$ is prime and
$\varepsilon \notin \mathbb{Q}$, the characteristic polynomial of
multiplication by $\varepsilon$ is its (irreducible) minimal
polynomial $X^3 - tX^2 + uX - 1$ (constant term $N(\varepsilon) = 1$
by Lemma~\ref{lem:normone}), with simple roots $\lambda, \mu, \bar\mu$.
Hence the cleared sequence
\[
z_k \;:=\; 3b\,c_2(\gamma\varepsilon^k) \;\in\; \mathbb{Z}
\]
satisfies the integer recurrence $z_{k+3} = t\,z_{k+2} - u\,z_{k+1} +
z_k$ with \emph{rational integer} coefficients --- the parameter
collapse that keeps every constant tame (in Sha's notation
\cite{Sha19}, the Galois-closure degree of the coefficient field is
$1$). The reversed sequence $v_n := z_{-n}$ satisfies the recurrence
with characteristic polynomial $X^3 - uX^2 + tX - 1$, the minimal
polynomial of $\varepsilon^{-1}$ --- \textbf{integral because
$N(\varepsilon) = 1$} --- with roots $\lambda^{-1}$ (of modulus
$e^{-R}$) and the conjugate pair $\mu^{-1}, \bar\mu^{-1}$ (of modulus
$e^{R/2}$).

%% ==================================================================
\section{The algorithm}\label{sec:algorithm}
%% ==================================================================

We display the algorithm, then discuss each step; the seven steps
follow the verified skeleton exactly. Lemmas A$'$, B, U, C and
Proposition F are stated and proved in Section~\ref{sec:lemmas}.

\subsection{Pseudocode}\label{sec:pseudocode}

\begin{verbatim}
Input:  cube-free d >= 2 (not a cube), m != 0, in binary;
        s = O(log d + log|m|).
Output: the complete list of (x, y) in Z^2 with x^3 - d y^3 = m.

Step 1  (Field data and infrastructure; no class group.)
        Factor d by trial division (cost O(d^{1/2}) poly = 2^{O(s)});
        write d = a b^2 with a, b coprime squarefree; extract the
        clearing constant 3b.
        Run BW88 Algorithm 2.13 on O_K: compute the Voronoi cycle of
        reduced principal ideals, the neighbor minima x_1, ..., x_p,
        the exact incremental products gamma_k = x_1 ... x_k, the
        fundamental unit eps = gamma_p, and the regulator R. [Lemma A']
        Expand eps in the integral basis (2^{O(s)} digits).

Step 2  (Norm-equation representatives.)
        Factor |m| by trial division. Enumerate all ideals I of O_K
        with N(I) = |m|  (at most d_3(|m|) of them; Dedekind basis,
        splitting at p | 3b read off from O_K/pO_K). For each I: test
        principality by reducing I and searching the cycle stock; if
        principal, extract the generator
        gamma_I = mu_I * gamma_{k-1}^{-1} and verify (gamma_I) = I by
        one exact HNF check.                             [Lemma A']

Step 3  (Unit reduction.)
        Replace each generator gamma_I by gamma' = gamma_I * eps^j with
        | log|sigma_1(gamma')| - (1/3)log|m| | <= R/2.   [Lemma U]

Step 4  (Recurrence data.)
        For each gamma', form z_k = 3b * c_2(gamma' eps^k): an integer
        sequence satisfying the order-3 recurrence with characteristic
        polynomial minpoly(eps); the reversed sequence v_n = z_{-n}
        satisfies the integral reversed recurrence (N(eps) = 1).
        Expand eps and z_0, z_1, z_2 to their 2^{O(s)} digits.

Step 5  (Forward window.)   k_+ := 3.                    [Proposition F]

Step 6  (Backward window.)  N_C := c (R + log|m| + log d + 1) *
                                   log(R + log|m| + log d + 2). [Lemma C]

Step 7  (The walk and the filter.)
        For each representative gamma' and each sign sigma in {+1, -1}:
        iterate the exact integer recurrences over k in [-N_C, k_+];
        whenever z_k = 0, write sigma * gamma' eps^k = x - y alpha and
        KEEP (x, y) if and only if x^3 - d y^3 = m exactly
        (mandatory norm-sign filter).
        Output the collected set of solutions.
\end{verbatim}

\subsection{Step 1 --- regulator and unit, no class group}
\label{sec:step1}

The step opens by factoring $d$ by trial division --- cost
$O(d^{1/2})\,\mathrm{poly}(s) = 2^{O(s)}$ --- which yields the
decomposition $d = ab^2$ and with it the clearing constant $3b$
consumed by Steps 2 and 4 (Section~\ref{sec:index}). At unit rank one
the class group is never needed: principality of
ideals is decided directly on the cycle of reduced principal ideals,
so no certification burden attaches to $\mathrm{Cl}_K$ at all. We
compute $R$ and the cycle by Buchmann--Williams Algorithm 2.13
\cite[Prop.~2.14]{BW88}: deterministic $O(R\,D^{\epsilon})$
bit-operations ($D = |\Delta|$), unconditional, published with proof.
By Section~\ref{sec:regulator}, $R = 2^{O(s)}$, so this is $2^{O(s)}$.
The fundamental unit in \emph{expanded} form costs $2^{O(s)}$ digits
--- affordable at this budget, and needed in Step 4 anyway. (Lenstra
\cite[Thm.~5.5]{Len92} --- whose deterministic exponent is
$|\Delta|^{3/4}$, not $1/2$, and which Lenstra himself hedges with
``appears to be true'' --- is a fallback citation only; the BW88 route
is the one we rely on.)

\subsection{Step 2 --- norm-equation representatives}\label{sec:step2}

Trial-divide $|m|$ (cost $2^{s/2}\mathrm{poly}(s)$). The ideals of
$O_K$ of norm $|m|$ number at most $d_3(|m|) = 2^{O(s)}$ (three-fold
divisor function; the counting argument is the proof line of
\cite[Thm.~6.5]{Len92}). They are enumerated prime-by-prime: for
$p \nmid 3b$, factor $x^3 - d \bmod p$; for $p \mid 3b$ (the primes
dividing the index), read the splitting off the finite algebra
$O_K/pO_K$ in the Dedekind basis; assemble prime-power blocks by
exponent vectors weighted by residue degrees, and multiply into
Hermite normal form, $\mathrm{poly}(s)$ per ideal. Per candidate
ideal, principality and a generator are obtained by the baby-step
cycle walk --- decision cost $O(R\,D^\epsilon)$ per candidate
\cite[Prop.~4.5]{BW88} --- with the tracked generator supplied by
Lemma A$'$. (The $R^{1/2}$-rate giant-step variant
\cite[Thm.~4.7]{BW88} is decision-only and is not used.)

\subsection{Step 3 --- unit reduction}\label{sec:step3}

Each generator is reduced along the one-dimensional log-unit lattice
(Lemma U, sharpened form): the reduced representative
$\gamma'$ satisfies
\[
\bigl|\log|\sigma_1(\gamma')| - \tfrac13\log|m|\bigr| \le \tfrac{R}{2},
\qquad\text{hence}\qquad
\bigl|\log|\sigma_2(\gamma')| - \tfrac13\log|m|\bigr| \le \tfrac{R}{4}
\ \ \text{and}\ \
\Bigl|\frac{\sigma_2(\gamma')}{\sigma_1(\gamma')}\Bigr| \le e^{3R/4}.
\]
The windows of Steps 5--6 apply to the \textbf{unit-reduced}
representatives --- this hypothesis is part of their statements.

\subsection{Step 4 --- orbit and recurrence}\label{sec:step4}

As in Section~\ref{sec:coordinate}: $z_k$ is the $\alpha^2$-coordinate
of $\gamma'(\pm\varepsilon)^k$, cleared by $3b$ (\textbf{not} by 3;
Section~\ref{sec:index}). The recurrence has order 3, characteristic
polynomial equal to the minimal polynomial of $\varepsilon$, simple
roots, and all Binet coefficients nonzero since $\gamma' \ne 0$ (DFT
weights of modulus $1/(3d^{2/3})$); the coefficients are rational
integers, so Sha's Galois parameter collapses to $1$.
$N(\varepsilon) = \lambda|\mu|^2 > 0$ forces $N(\varepsilon) = +1$
(Lemma~\ref{lem:normone}), which makes the reversed recurrence
integral. Expanding $\varepsilon$ and $z_0, z_1, z_2$ to their
$2^{O(s)}$ digits is part of this step's stated cost.

\subsection{Steps 5 and 6 --- the two windows}\label{sec:steps56}

The forward window is elementary and absolute: vanishing at $k > 0$
forces $\lambda^{3k/2} \le 2e^{3R/4}$, giving $k_+ \le 3$
(Proposition F --- which is exactly the empirical maximum,
Section~\ref{sec:prototype}). The backward window is Lemma C:
$v_n \ne 0$ for all
$n > N_C = c\,(R + \log|m| + \log d + 1)\log(R + \log|m| + \log d +
2)$, an absolute explicitly computable $c$ --- quasi-linear in $R$,
hence $N_C = \widetilde O(R + s) = 2^{O(s)}$.

\subsection{Step 7 --- the walk}\label{sec:step7}

For every representative $\times$ sign $\times$ index
$k \in [-N_C, k_+]$: one exact integer recurrence iteration
($2^{O(s)}$ digits each, $2^{O(s)}$ total), collect the hits
$z_k = 0$, recover $(x, y)$ from the $1$- and $\alpha$-coordinates,
and apply the \textbf{mandatory final filter} $x^3 - dy^3 = m$: the
walk also emits norm-$(-m)$ elements, and only the filter separates
the two signs. (Since $c_2$ is linear, the $\pm$ sequences share
their zero set; the sign branch matters at recovery, where
$\gamma \mapsto -\gamma$ swaps the norms $\pm m$.)

Completeness is ideal-theoretic and library-free: any solution
$\gamma = x - y\alpha \in \mathbb{Z}[\alpha] \subseteq O_K$ generates
an ideal of norm $|m|$ appearing in Step 2's list; the unit orbits
under $\langle -1, \varepsilon\rangle$ --- the full unit group,
torsion $\{\pm1\}$ by the real embedding --- cover all generators of
each principal ideal; and $\gamma \mapsto -\gamma$ swaps $\pm m$,
which the $\pm$-walk covers. The formal proof is
Section~\ref{sec:proofmain}.

%% ==================================================================
\section{The supporting lemmas}\label{sec:lemmas}
%% ==================================================================

Throughout this section $\gamma$ denotes a nonzero element of $O_K$
with $N(\gamma) = \pm m$, and $b_1, b_2, \bar b_2$ the ($3b$-cleared)
Binet coefficients of its orbit sequence
(Section~\ref{sec:coordinate}), all nonzero.

\subsection{Lemma B: root-of-unity exclusion}\label{sec:lemB}

\begin{lemmaB}
The quotient $\bar\mu/\mu$ of the complex embeddings of $\varepsilon$
is not a root of unity.
\end{lemmaB}

\begin{proof}
Suppose $(\bar\mu/\mu)^t = 1$ for some $t \ge 1$. Then
$\sigma_2(\varepsilon^t) = \sigma_3(\varepsilon^t) =
\overline{\sigma_2(\varepsilon^t)}$, so $\sigma_2(\varepsilon^t)$ is
real. (The contradiction target is ``$\sigma_2(\varepsilon^t)$ real''
--- not ``$\varepsilon^t$ real'': $\sigma_1(\varepsilon^t)$ is always
real, so that statement would be empty.) Since $[K:\mathbb{Q}] = 3$
is prime, either $\varepsilon^t \in \mathbb{Q}$, in which case
$\varepsilon^t$ is a rational unit, $\varepsilon^t = \pm 1$,
contradicting $|\sigma_1(\varepsilon^t)| = e^{tR} > 1$; or
$\mathbb{Q}(\varepsilon^t) = K$, in which case the three conjugates
of $\varepsilon^t$ are pairwise distinct, contradicting
$\sigma_2(\varepsilon^t) = \sigma_3(\varepsilon^t)$.
\end{proof}

The same argument applied to $|\lambda/\mu| = e^{3R/2} \ne 1$ shows
the full sequence is non-degenerate (no ratio of distinct roots is a
root of unity), so its zero set is finite; but the proof below needs
only the $\bar\mu/\mu$ case.

\subsection{Lemma U: unit reduction at rank one}\label{sec:lemU}

\begin{lemmaU}
Let $K$ be a complex cubic field, $\varepsilon$ normalized with
$\lambda = \sigma_1(\varepsilon) > 1$, $R = \log\lambda$. For every
$\gamma \in O_K$ with $N(\gamma) = m' \ne 0$ there is
$j \in \mathbb{Z}$ such that $\gamma' = \gamma\varepsilon^j$ satisfies
\[
\bigl|\log|\sigma_1(\gamma')| - \tfrac13\log|m'|\bigr| \;\le\;
\tfrac{R}{2}.
\]
Consequently
$\bigl|\log|\sigma_2(\gamma')| - \tfrac13\log|m'|\bigr| \le R/4$,
$\ |\sigma_2(\gamma')/\sigma_1(\gamma')| \le e^{3R/4}$, and
$h(\gamma') \le \tfrac13\log|m'| + R + O(1)$.
\end{lemmaU}

\begin{proof}
$|\sigma_1(\gamma)|\,|\sigma_2(\gamma)|^2 = |m'|$, and for
$t(\gamma) := \log|\sigma_1(\gamma)| - \tfrac13\log|m'|$ we have
$t(\gamma\varepsilon^j) = t(\gamma) + jR$; take
$j = -\mathrm{round}(t(\gamma)/R)$. The consequences follow from
$|\sigma_2|^2 = |m'|/|\sigma_1|$ and from bounding the house of
$\gamma'$.
\end{proof}

This is a one-dimensional reduction in the log-unit lattice ---
genuinely trivial as mathematics. Two computational remarks are part
of the record: computing $j$ rigorously requires $\sigma_1(\gamma)$
to polynomially many digits via interval arithmetic, cost $2^{O(s)}$
on $2^{O(s)}$-digit inputs; and a borderline rounding at
half-integers only widens the forward window by 1 --- a slack already
accounted for in Proposition F.

\subsection{Lemma A$'$: tracked generators along the Voronoi cycle}
\label{sec:lemAprime}

\begin{lemmaAprime}
Fix the degree and unit rank one, and let $p$ be the cycle length, so
$p < 2R/\log\tau + 1$ with $\tau$ the golden ratio (Williams:
$\varepsilon_0 > \tau^{p/2}$). The neighbor minima $x_1, \ldots, x_p$
of the reduced-ideal cycle each carry $\mathrm{poly}(s)$-size data
\cite[Prop.~2.11]{BW88}, and the incremental exact products
$\gamma_k = x_1\cdots x_k$ ($2^{O(s)}$ digits each) satisfy: the
$(k{+}1)$-st reduced principal ideal is $I_{k+1} = (\gamma_k)^{-1}$,
and $\varepsilon = \gamma_p$. Hence every principality witness and
$\varepsilon$ itself are computable, expanded, in deterministic total
time $2^{O(s)}$, with exact HNF verification of each ideal identity.
\end{lemmaAprime}

\begin{proof}[Proof (with the corrected orientation and
certification)]
The Voronoi chain at rank one is $I_1 = O_K$ and
$I_{k+1} = (1/x_k)\,I_k$, where $x_k$ is the neighbor minimum of
$I_k$, \textbf{computed exactly by BW88 Algorithm 2.13 / Williams
\cite{Wil85}}. Exactness of the neighbor step is what certifies
completeness of the stock: floats used for steering alone could skip
a minimum and still pass every HNF test, so the per-step computation
must be the exact one. Telescoping, $\gamma_k = x_1\cdots x_k$ is the
$(k{+}1)$-st minimum and $I_{k+1} = (\gamma_k)^{-1}$; after a full
cycle, $\gamma_p = \varepsilon$ exactly, under the $\sigma_1 > 0$,
increasing-$\lambda$ normalization. $\mathrm{poly}(s)$-size of each
$x_k$ follows from \cite[Prop.~2.4]{BW88}
($|N(x)| \le \sqrt{D}\,N(\mathfrak{a})$), \cite[Prop.~2.7(i)]{BW88}
(log-step $< \log\sqrt{D}$), and denominators $\le \sqrt{D}$. Digit
growth of the exact products is additive,
$O(p\cdot\mathrm{poly}(s)) = 2^{O(s)}$ per $\gamma_k$ and $2^{O(s)}$
in total, including the per-step exact HNF chain identity
$I_{k+1}\cdot(x_k) = I_k$.

\textbf{Candidate-side recipe} (Step 2's need): given an ideal $I$ of
norm $|m|$, reduce exactly to obtain the relative minimum $\mu_I$
($\mathrm{poly}(s)$-size); $I$ is principal if and only if the
reduction lands on the cycle stock at some position $k$, and then its
generator is
\[
g \;=\; \mu_I\cdot\gamma_{k-1}^{-1}
\]
(exact division on $2^{O(s)}$-digit integers), certified by one
terminal exact HNF check $(g) = I$.
\end{proof}

Three remarks on the constants and the design, verified against the
primary sources.

\begin{enumerate}
\item \textbf{Cycle length.} \cite[Cor.~2.8]{BW88} gives
  $R/\log\sqrt{D} < p \le jR/c_1$ with $(j, c_1) = (7, \log 4)$ at
  degree 3 --- constants confirmed against the Math.\ Comp.\ text and
  against Williams \cite{Wil86}, whose scope is \emph{all} complex
  cubic fields and arbitrary reduced lattices (the paper also has the
  sharper $\theta_8 > 4$, and $\varepsilon_0 > \tau^{p/2}$, whence
  $p < 2R/\log\tau + 1 \approx 4.16R$ as used in the statement). A
  venue correction carried in our records: Williams' \emph{Continued
  fractions and number-theoretic computations} is Rocky Mountain J.
  Math.\ 15 (1985) \cite{Wil85}, not Pacific J. Math.; the Pacific J.
  Math.\ paper is \cite{Wil86}.
\item \textbf{Why plain products.} An earlier plan routed height
  bounds through composition-then-reduce distance-slip estimates. The
  simplification adopted here: at the $2^{O(s)}$ budget
  one simply takes the plain incremental exact products of the at
  most $7R/\log 4$ neighbor minima along the cycle. Compact
  representations (Thiel) are a luxury needed only for
  $\mathrm{poly}$-size \emph{certificates}, not for the EXP
  algorithm.
\item \textbf{What Lemma A$'$ delivers to Step 1.} The full cycle
  walk computes $\varepsilon = \gamma_p$ and $R$ along the way; Steps
  1 and 2 share one infrastructure computation.
\end{enumerate}

\subsection{Proposition F: the forward window is $O(1)$}
\label{sec:propF}

\begin{propositionF}
Let $\gamma'$ be a unit-reduced representative (Lemma U) and $z_k$
the $3b$-cleared $\alpha^2$-coordinate sequence of
$\pm\gamma'\varepsilon^k$. If $z_k = 0$ with $k > 0$, then
\[
\lambda^{3k/2} \;\le\;
2\,\Bigl|\frac{\sigma_2(\gamma')}{\sigma_1(\gamma')}\Bigr|
\;\le\; 2e^{3R/4},
\qquad\text{so}\qquad
k \;\le\; \frac12 + \frac{2\log 2}{3R} \;\le\; 2
\]
using the sharp complex-cubic floor $R \ge 0.2811995743$
\cite{ADF16}. Allowing the one-step rounding slack of Lemma U:
\[
k_+ \;\le\; 3,
\]
which is exactly the maximum observed empirically
(Section~\ref{sec:prototype}).
\end{propositionF}

\begin{proof}
By the DFT form (Section~\ref{sec:coordinate}), $z_k = 0$ means
$b_1\lambda^k = -(b_2\mu^k + \bar b_2\bar\mu^k)$, whence
$|\sigma_1(\gamma')|\lambda^k \le 2|\sigma_2(\gamma')|\lambda^{-k/2}$
(the weights share the modulus $1/(3d^{2/3})$, which cancels; the
$3b$ clearing also cancels). Rearranged: $\lambda^{3k/2} \le
2|\sigma_2(\gamma')/\sigma_1(\gamma')| \le 2e^{3R/4}$ by Lemma U.
Taking logarithms, $k \le \tfrac12 + \tfrac{2\log 2}{3R}$, and with
$R \ge 0.28119$ the right side is $\le 2.15$, so $k \le 2$; the
possible half-integer rounding in Lemma U widens the admissible
window by at most one index.
\end{proof}

Note the phenomenon the empirical data confirm: after unit reduction
the $R$'s cancel --- the forward window
is not merely $\mathrm{poly}(s)$, it is \emph{bounded by an absolute
constant}.

\subsection{Lemma C: the backward window, quasi-linear in $R$}
\label{sec:lemC}

\begin{lemmaC}
Let $\gamma$ be a unit-reduced representative with
$N(\gamma) = \pm m$, and let $v_n$ be the $3b$-cleared
$\alpha^2$-coordinate sequence of $\gamma\varepsilon^{-n}$
($n \ge 0$). There is an absolute, explicitly computable constant $c$
such that $v_n \ne 0$ for all
\[
n \;>\; N_C \;:=\; c\,\bigl(R + \log|m| + \log d + 1\bigr)\,
\log\bigl(R + \log|m| + \log d + 2\bigr).
\]
\end{lemmaC}

\begin{proof}
Write $v_n = b_1\lambda^{-n} + b_2\mu^{-n} + \bar b_2\bar\mu^{-n}$
with all $b_j \ne 0$ ($\gamma \ne 0$;
Section~\ref{sec:coordinate}), and set
$\mu^{-1} = e^{R/2}e^{i\theta}$, $b_2 = |b_2|e^{i\varphi}$. The two
maximal-modulus roots of the reversed recurrence are the conjugate
pair $\mu^{-1}, \bar\mu^{-1}$ (modulus $e^{R/2}$), the third root
$\lambda^{-1}$ is smaller by the exact factor $e^{3R/2}$, and the
quotient $\bar\mu/\mu$ of the maximal pair is not a root of unity
(Lemma B).

If $v_n = 0$, then
$2|b_2|e^{nR/2}|\cos(\varphi + n\theta)| = |b_1|e^{-nR}$, so with
$|\cos x| \ge \tfrac{2}{\pi}\,\mathrm{dist}(x, \tfrac{\pi}{2} +
\pi\mathbb{Z})$:
\begin{equation}
\mathrm{dist}\Bigl(\varphi + n\theta,\ \tfrac{\pi}{2} +
\pi\mathbb{Z}\Bigr)
\;\le\; \frac{\pi}{4}\,\frac{|b_1|}{|b_2|}\,e^{-3nR/2}.
\tag{$\dagger$}
\end{equation}

The distance on the left is measured exactly by the linear form in
logarithms (with the corrected conjugation and exact factor 2, both
verified numerically on two fields):
\[
\Lambda_n \;=\; a\,\log(-1) \;+\; \log\!\frac{b_2}{\bar b_2}
\;+\; n\,\log\!\frac{\bar\mu}{\mu},
\qquad a \in \mathbb{Z} \text{ odd},\ |a| \le n + 2,
\]
with
\[
|\Lambda_n| \;=\; 2\,\mathrm{dist}\Bigl(\varphi + n\theta,\
\tfrac{\pi}{2} + \pi\mathbb{Z}\Bigr)
\quad\text{exactly.}
\]

\textbf{The degenerate case is void.} $\Lambda_n = 0$ if and only if
$\cos(\varphi + n\theta) = 0$, i.e.\ the conjugate-pair term of $v_n$
vanishes --- but then $v_n = b_1\lambda^{-n} \ne 0$. So at any actual
zero $v_n = 0$ we always have $\Lambda_n \ne 0$, and Matveev's
theorem applies.

Apply Matveev \cite{Mat00} (as restated in \cite[\S 2.4]{Sha19}) with
$k = 3$ logarithms in the field
$L = \mathbb{Q}(\sqrt[3]{d}, \omega)$ of degree
\[
D_0 = 6.
\]
$L$ \textbf{is} the Galois closure of $K$: $\sqrt{\mathrm{disc}(K)}$
generates $\mathbb{Q}(\sqrt{-3}) = \mathbb{Q}(\omega)$ for
\emph{every} pure cubic field (the discriminant is $-3$ times a
square in both Dedekind types), and $L$ is stable under complex
conjugation and contains $-1$, $b_2/\bar b_2$, and $\bar\mu/\mu$. The
conclusion is
\[
\log|\Lambda_n| \;>\; -\,C_0\,A_1A_2A_3\,\log\bigl(e(n+2)\bigr),
\qquad C_0 = 2^{40}\,D_0^2\,\log(eD_0),
\]
where the $2^{40}$-versus-$2^{38}$ slack absorbs the factor 2 in
$|\Lambda_n| = 2\,\mathrm{dist}$, and:

\begin{itemize}
\item $A_1 = \pi$;
\item $A_2 \le D_0\,h(b_2/\bar b_2) + \pi \le 12H_1 + \pi$, where
  (heights are Galois-invariant, and the $3b$-clearing adds at most
  $\log 3d$):
  \[
  h(b_2) \;\le\; h(\gamma) + \log 3 + \tfrac23\log d + \log 3d
  \;=:\; H_1 \;=\; O\bigl(h(\gamma) + \log d + 1\bigr);
  \]
\item $A_3 = \max\bigl(D_0\,h(\bar\mu/\mu),\
  |\log(\bar\mu/\mu)|\bigr) \le \max(4R, \pi) = 4R$, using
  $h(\bar\mu/\mu) \le 2h(\varepsilon) = 2R/3$
  ($h(\varepsilon) = R/3$; the complex place carries weight 2).
  Matveev's side condition $A_3 \ge |\log(\bar\mu/\mu)|$ --- a
  quantity which can approach $\pi$ --- is met because pure cubic
  fields have $|\mathrm{disc}| \ge 108$, hence $R > 0.79$ (the
  exceptional fields of \cite{ADF16} at discriminants $-23, -31,
  -44$ are not pure), so $4R \ge 3.16 > \pi$ ---
  \textbf{margin $0.018$}.
\end{itemize}

Combining Matveev's lower bound with $(\dagger)$ and
$|\log(|b_1|/|b_2|)| \le 12H_1$ (via
$|\log|b|| \le \deg(b)\cdot h(b)$ and
$\deg \le 6$):\footnote{The displayed combined inequality retains the
  verified notes' doubled parameters $24H_1 + \pi$ and $8R$, a
  factor-4 safety slack over the minimal admissible Matveev choices
  $A_2 = 12H_1 + \pi$, $A_3 = 4R$ listed above; the slack is absorbed
  into the final constant evaluation flagged in
  Remark~\ref{rem:cchain} and does not affect the shape of $N_C$.}
\[
\frac{3R}{2}\,n \;\le\; 12H_1 + \log\frac{\pi}{4}
\;+\; C_0\,\pi\,(24H_1 + \pi)\,(8R)\,\log\bigl(e(n+2)\bigr).
\]
Divide by $3R/2$: \textbf{the $R$ cancels in the Matveev term},
leaving
\[
n \;\le\; \frac{8H_1}{R} \;+\; c_2\,(H_1 + 1)\,
\log\bigl(e(n+2)\bigr),
\qquad c_2 \text{ absolute and explicit.}
\]
Resolving $n$ against $\log n$ (via $x \ge 2A\log A \Rightarrow
x/\log x \ge A$; the term $8H_1/R$ is $O(H_1)$ by the regulator
floor) and inserting the unit-reduction bound
$H_1 = O(\log|m| + R + \log d)$ (Lemma U) gives $N_C$ as stated. In
particular
\[
N_C \;=\; \widetilde O(R + s) \;=\; 2^{O(s)}
\]
--- quasi-linear in $R$, in place of the fatal $e^{22.5R}$ of the
black-box route (Remark~\ref{rem:blackbox}).
\end{proof}

\emph{Final write-up bookkeeping, carried forward visibly from the
verified notes:} the absolute constant $c$ is traced through the
chain $2^{40}\cdot 36\cdot\log(6e)\cdot\pi\cdot(12H_1+\pi)\cdot 4R$
(before the $R$-cancellation), with $D_0 = 6$ throughout, and the
branch conventions for $\log(-1)$ absorbed into $|a| \le n + 2$ ---
all mechanical; the numeric evaluation is carried out in
Remark~\ref{rem:cchain}.

\subsection{Remark: why the black-box window does not suffice (a
  preserved lesson)}\label{rem:blackbox}

An earlier draft of this work cited Sha's Theorem 1.2 \cite{Sha19} as
a black box, with ``window $2^{O(s)}$''. The claim was wrong in the
way that matters most in this subject: the $s$ in that statement was
the extraction's \emph{redefined} size
$\max(R, h(\gamma), \log|m|, \log d) \supseteq R$, and
$R = 2^{\Theta(s_{\mathrm{input}})}$. In honest input size, the
black-box window $e^{22.5R}$ is \textbf{doubly exponential} --- a
measured instance: $d = 9986$ has $R \approx 2605$ and a black-box
window exceeding $10^{25{,}454}$ steps. The exponential loss is
entirely an artifact of the worst-case Mahler-type root-separation
quantity in Sha's theorems, which is $e^{-\Theta(R)}$ for our
instances even though the \emph{true} ratio gap is $3R/2$ ---
exponentially larger. Substituting the true gap into Sha's own
endgame (a change of the proofs, not of the theorems --- which is why
Lemma C is stated as ours, with Sha's Lemma 2.8 and the Matveev
application as extracted ingredients) is what produces the
quasi-linear window. This was the third instance of the same trap in
this program's history, and it has earned a name we record for the
reader: \textbf{always ask exponential in what.}

%% ==================================================================
\section{Proof of the Main Theorem}\label{sec:proofmain}
%% ==================================================================

\begin{theoremmain}
There is a deterministic, unconditional algorithm that, given
cube-free $d \ge 2$ (not a cube) and $m \ne 0$ in binary, with
$s = O(\log d + \log|m|)$, decides whether $x^3 - dy^3 = m$ has a
solution in $\mathbb{Z}^2$ --- listing all solutions --- in time
$2^{O(s)}$.
\end{theoremmain}

\begin{proof}
We show the algorithm of Section~\ref{sec:algorithm} is sound,
complete, and runs in time $2^{O(s)}$.

\textbf{Soundness.} Every pair the algorithm outputs has passed the
exact integer filter $x^3 - dy^3 = m$ in Step 7, so no false
positives occur --- in particular the norm-$(-m)$ elements that the
walk necessarily emits are discarded there.

\textbf{Completeness (ideal-theoretic).} Let
$(x, y) \in \mathbb{Z}^2$ satisfy $x^3 - dy^3 = m$, and set
$\gamma = x - y\alpha \in \mathbb{Z}[\alpha] \subseteq O_K$, so
$N(\gamma) = m$. The ideal $(\gamma)$ has norm $|m|$ and therefore
appears in Step 2's enumeration; it is principal, so Step 2
(Lemma A$'$) finds a generator $\gamma_I$ with
$(\gamma_I) = (\gamma)$, and Step 3 replaces it by the unit-reduced
$\gamma'$ in the same orbit. Two generators of the same ideal differ
by a unit, and the unit group is $\langle -1, \varepsilon\rangle$ ---
the full unit group, since the torsion is $\{\pm 1\}$ by the real
embedding --- so
\[
\gamma \;=\; \pm\,\gamma'\,\varepsilon^{k}
\quad\text{for some } k \in \mathbb{Z}.
\]
The $\alpha^2$-coordinate of $\gamma$ vanishes (Thue shape), so the
cleared sequence of $\gamma'$ has $z_k = 0$. Since $\gamma'$ is
unit-reduced, Proposition F bounds the positive side,
$k \le k_+ = 3$, and Lemma C bounds the negative side,
$-k \le N_C$; both windows apply precisely because they are stated
for unit-reduced representatives. Hence $k \in [-N_C, k_+]$, the walk
of Step 7 visits index $k$, detects $z_k = 0$, and recovers from
$\pm\gamma'\varepsilon^k$ both candidate pairs $\pm(x, y)$; the
filter retains exactly the one(s) of norm $m$ --- among them
$(x, y)$. Both norm signs are covered: if $N(\gamma') = -m$, the
solution appears through the sign branch
$\gamma = -\gamma'\varepsilon^k$, since $\gamma \mapsto -\gamma$
swaps $\pm m$ at odd degree. No step of this argument appeals to any
external solver; completeness is purely ideal-theoretic.

\textbf{Cost accounting.} Write $D = |\Delta| \le 27d^2 = 2^{O(s)}$
and recall $R \le hR = 2^{O(s)}$ (Section~\ref{sec:regulator}).

\begin{itemize}
\item \emph{Step 1:} trial-division factorization of $d$,
  $O(d^{1/2})\,\mathrm{poly}(s) = 2^{O(s)}$; then
  $O(R\,D^\epsilon) = 2^{O(s)}$ bit-operations
  \cite[Alg.~2.13, Prop.~2.14]{BW88}; expanding $\varepsilon$ costs
  $2^{O(s)}$ digits (Lemma A$'$).
\item \emph{Step 2:} trial division $2^{s/2}\mathrm{poly}(s)$; at
  most $d_3(|m|) = 2^{O(s)}$ candidate ideals, each assembled in
  $\mathrm{poly}(s)$; per candidate one cycle search,
  $O(R\,D^\epsilon) = 2^{O(s)}$ \cite[Prop.~4.5]{BW88}, with
  generator extraction and one exact HNF verification (Lemma A$'$).
\item \emph{Step 3:} interval arithmetic to polynomially many digits
  on $2^{O(s)}$-digit inputs: $2^{O(s)}$ (Lemma U).
\item \emph{Step 4:} expanding $z_0, z_1, z_2$ and the recurrence
  coefficients: $2^{O(s)}$ digits, stated cost of the step.
\item \emph{Steps 5--6:} window sizes $k_+ \le 3$ and
  $N_C = \widetilde O(R + s) = 2^{O(s)}$.
\item \emph{Step 7:} at most $(N_C + k_+ + 1)$ iterations per
  representative and sign; each iteration is exact integer
  arithmetic on $2^{O(s)}$-digit numbers; total
  $2^{O(s)}\cdot 2^{O(s)} = 2^{O(s)}$.
\end{itemize}

The product of a $2^{O(s)}$ number of representatives with
$2^{O(s)}$ work each is $2^{O(s)}$, completing the proof.
\end{proof}

We record the honest caveat once more: to \emph{run} Step 6 one needs
a numeric value for $c$; the constant chain of Lemma C supplies one
mechanically, and its final evaluation is deferred
(Remark~\ref{rem:cchain}). The theorem is unconditional and its proof
complete; the window constant is semi-explicit --- explicit modulo
final constant evaluation.

%% ==================================================================
\section{Computational companion}\label{sec:companion}
%% ==================================================================

Two experiments and one independent reimplementation accompany the
proof. Neither is part of the logical argument; both
shaped it.

\subsection{The orbit-walk prototype: 1120/1120}\label{sec:prototype}

\texttt{experiments/cubic\_orbit\_prototype.py} implements the orbit
walk with PARI/GP \cite{PARI} supplying the algebraic data:
representatives from \texttt{bnfisintnorm} over
\texttt{bnfinit(x\^{}3 - d, 1)}, the $\pm\varepsilon^{\pm k}$ walk
over $k \in [-60, 60]$, Thue-shape detection on the
$\alpha^2$-coordinate, and the exact final filter $x^3 - dy^3 = m$.
The validation grid is
\[
d \in \{2, 3, 5, 6, 7, 10, 11, 12, 13, 15, 17, 19, 20, 22\},
\qquad m \in [-40, 40] \setminus \{0\}
\]
--- 14 fields $\times$ 80 right-hand sides $=$ \textbf{1120 cases}
--- and the reference is PARI's \texttt{thue} over
\texttt{thueinit(x\^{}3 - d, 1)}, whose flag-1 certification is
unconditional. Result: \textbf{1120/1120 exact solution-set
agreement}, and across all hits the maximum vanishing index was
\[
\max |k| = 3
\]
--- exactly the bound $k_+ \le 3$ that Proposition F later proved
(the prototype's fixed cutoff $K_0 = 60$ was the honest gap between
prototype and theorem; the two windows of Section~\ref{sec:lemmas}
close it). We stress the boundary between the proven algorithm and
this practical code: the cutoff $K_0 = 60$ operates under the strong
empirical window conjecture of Remark~\ref{rem:smale}, not under the
proven bound of Lemma C; the proven window $N_C$ --- of order
$10^{17}\cdot H_1\log H_1$ steps once the constants of
Remark~\ref{rem:cchain} are inserted --- is $2^{O(s)}$ in theory but
is not what the prototype executes. For completeness:
\texttt{bnfisintnorm} returns a complete
system modulo units of \emph{positive} norm, and the
$\langle -1, \varepsilon \rangle$ walk covers a superset of those
orbits.

\subsection{Unit growth: the wall being dodged}\label{sec:unitgrowth}

\texttt{experiments/thue\_unit\_growth.py} charts why any method that
must expand unit data pays exponentially: for the family
$x^3 - dy^3 = 1$ it records $d$ whenever the digit-size of the
fundamental unit's coefficients sets a record (regulators from
\texttt{bnfinit}, GRH-conditional; every record row re-certified
unconditionally with \texttt{bnfcertify}). By dense size $s = 20$ (at
$d = 1721$, $R \approx 3669.4$) the coefficients reach \textbf{1593
digits}, all records certified --- a steeper envelope than the Pell
layer (roughly 145 digits at the same size), one rung down. The data
are in \texttt{experiments/data/cubic\_unit\_records.csv}. This is
precisely the obstruction the infrastructure route dodges: the
algorithm expands $\varepsilon$ only inside its stated $2^{O(s)}$
budget, and every per-step object it manipulates is polynomially
sized (Lemma A$'$).

\subsection{An independent reimplementation: 13/13}
\label{sec:reimpl}

An independent \textbf{from-scratch mini-implementation} of the
algorithm, written separately from the prototype, was validated
against certified \texttt{thue}: \textbf{13/13 exact agreements},
including the high-index fields
\[
d = 12,\ 45,\ 175 \qquad
([\,O_K : \mathbb{Z}[\alpha]\,] = 2,\ 3,\ 5),
\]
precisely the fields at which the naive ``clear by 3'' recipe fails
(Section~\ref{sec:index}) --- so the validation deliberately covers
the corrected code path. The same reimplementation verified the
linear-form analysis of Lemma C numerically on two fields and
validated the Binet integer sequence against a Thue-shape zero.

%% ==================================================================
\section{Remarks and open problems}\label{sec:remarks}
%% ==================================================================

\subsection{General complex cubic forms (rank 1)}\label{rem:general}

The natural next target is the general irreducible cubic Thue
equation of negative discriminant, $F(x, y) = m$: the associated
cubic field is again complex, the unit rank is again 1, and the
entire skeleton --- cycle walk, unit reduction, two-sided window,
exact walk --- carries over. What grows is bookkeeping: the ring
$\mathbb{Z}[\alpha]$ attached to a general form is no longer
generated by a pure cube root, the index and clearing analysis of
Section~\ref{sec:index} must be redone for the form's ring, and the
height accounting in $H_1$ acquires the form's coefficients. We see
no conceptual obstruction, and state it as the immediate open
problem: extend Theorem~\ref{thm:main} to all complex cubic forms.

\subsection{Totally real cubics (rank 2) are the honest hard case}
\label{rem:totallyreal}

For totally real cubic forms the unit rank is 2 and the picture
changes qualitatively: the vanishing locus lives on a
$\mathbb{Z}^2$-lattice of units, the walk has no canonical successor,
and no dominant root separates a forward window --- the problem takes
on the flavor of S-unit equations rather than of a one-dimensional
Skolem instance. We know of no route to a $2^{O(s)}$ bound there and
explicitly defer it; nothing in this paper's method addresses it.

\subsection{The constant chain (evaluated)}\label{rem:cchain}

The constants of Lemma C are now fully explicit; we record the
evaluation. With $D_0 = 6$ (exact:
$\mathbb{Q}(\sqrt[3]{d}, \omega)$ \emph{is} the Galois closure,
Section~\ref{sec:lemC}), the absolute constant $c$ in $N_C$ is traced
through
\[
2^{40}\cdot 36\cdot\log(6e)\cdot\pi\cdot(12H_1 + \pi)\cdot 4R
\]
--- i.e.\ $C_0 = 2^{40}D_0^2\log(eD_0) = 2^{40}\cdot 36\,(1 + \log 6)
\approx 1.105\times 10^{14}$ against $A_1 = \pi$,
$A_2 = 12H_1 + \pi$, $A_3 = 4R$ --- before the $R$-cancellation.
Dividing the zero inequality
$\tfrac{3R}{2}\,n \le 12H_1 + \log\tfrac{\pi}{4} + C_0\,\pi\,(12H_1 +
\pi)(4R)\log(e(n+2))$ by $3R/2$ cancels the $R$ in the Matveev term
and leaves
\[
n \;\le\; \frac{8H_1}{R} \;+\; K\,(12H_1 + \pi)\,
\log\bigl(e(n+2)\bigr),
\qquad K \;:=\; \tfrac{2}{3}\,C_0\,\pi\cdot 4
\;=\; 2^{40}\cdot 96\,\pi\,(1 + \log 6)
\;\approx\; 9.258\times 10^{14}.
\]
Resolving $n$ against $\log n$ exactly as in the proof of Lemma C
(via $x \ge 2A\log A \Rightarrow x/\log x \ge A$; the term $8H_1/R$
is $O(H_1)$ by the regulator floor) yields the explicit window
\[
N_C \;\le\; \max\Bigl(\frac{16H_1}{R},\
4K\,(12H_1 + \pi)\,\log\bigl(2eK\,(12H_1 + \pi)\bigr)\Bigr),
\]
whose leading coefficient $48K \approx 4.44\times 10^{16}$ multiplies
$H_1\log(\cdot)$; under the factor-4 safety slack of the displayed
combined inequality in the proof of Lemma C (parameters
$24H_1 + \pi$ and $8R$), $K$ is replaced by
$4K \approx 3.703\times 10^{15}$ and the leading coefficient by
$\approx 1.78\times 10^{17}$. Either way the coefficient governing
$N_C$ is of order $10^{17}$, and the numeric value of $c$ is thereby
fixed; the elementary constants of Lemma U and Proposition F and the
branch conventions for $\log(-1)$ remain absorbed into
$|a| \le n+2$. Sha remarks \cite[\S 2.4]{Sha19} that sharper
three-logarithm bounds (Mignotte's kit) carry side conditions that
need not hold, so no off-the-shelf improvement to the $2^{40}$ is
available.

\subsection{Relation to Smale's fifth problem}\label{rem:smale}

In the stratification of Smale's fifth problem that this program
maintains, genus-0 components with at least three places at infinity
--- the stratum that classically reduces to Thue equations --- form
the first wall past the conic/Pell layer. Theorem~\ref{thm:main}
breaches that wall at its thinnest point: one fixed family, degree 3,
pure. The strong empirical window conjecture --- every Thue-shape
index of a unit-reduced representative satisfies
$|k| \le C(1 + \log|m|/R)$ --- remains open (max observed index: 3),
as do the general-form extensions above, and, a fortiori, the uniform
question of the full stratum: is
$\{f : \text{every component of genus } 0\}$ decidable in time
$2^{O(s)}$ uniformly in the degree? The order-3 Skolem instance at
the heart of our walk is classically decidable
(Mignotte--Shorey--Tijdeman \cite{MST84}, Vereshchagin \cite{Ver85});
our contribution is not decidability but an instance-specific, fully
explicit, two-sided window that is exponential in the \emph{input},
not in the regulator. We have not pursued poly-size certificates;
compact representations are exactly the extra ingredient a
certificate version would need (Lemma A$'$), a further direction we
leave open.

\subsection{Beyond exponential time: the named obstructions}
\label{rem:beyond}

Deterministic \emph{polynomial-time} decision for pure cubic Thue
equations is not a rewrite of this paper away: it faces three
independent open problems, which we name so that the reader can see
exactly where the exponential lives.

\begin{enumerate}
\item \textbf{The unit.} Polynomial-time computation of the regulator
  and the fundamental unit --- even in compact representation --- is
  equivalent to the principal ideal problem and is a major open
  problem of computational number theory, classically believed hard:
  it underlies infrastructure-based cryptography, and polynomial time
  is known only quantumly, by Hallgren-type algorithms \cite{Hal07}.
  Note that even negative-Pell \emph{decision}, one degree down, is in
  $\mathrm{NP} \cap \mathrm{coNP}$ (Lagarias \cite{Lag79}) with no
  classical polynomial algorithm known.
\item \textbf{The zero test.} Certified zero-testing of the recurrence
  values against their $2^{O(s)}$-digit height bounds via modular
  evaluation needs exponentially many primes; per-prime testing
  yields only randomized decisions (coRP-style, as in
  straight-line-program zero-testing), whose derandomization is open.
\item \textbf{The window.} The scan window is quasi-linear in
  $R = 2^{\Theta(s)}$ unless the strong window conjecture of
  Remark~\ref{rem:smale} --- empirical maximum index 3,
  Section~\ref{sec:prototype} --- is proven.
\end{enumerate}

The natural next target is therefore NP membership --- a cubic
analogue of Lagarias's Pell certificates \cite{Lag79} --- for which
the strong window conjecture is the enabler.

%% ==================================================================
%% References
%% ==================================================================

\begin{thebibliography}{TdW89}

\bibitem[ADF16]{ADF16} S. Astudillo, F. D\'iaz y D\'iaz, E. Friedman,
  \emph{Sharp lower bounds for regulators of small-degree number
  fields}, J. Number Theory 167 (2016) 232--258. (Theorem 7 and
  Table 1: the two regulator floors of
  Section~\ref{sec:regulator}.)

\bibitem[AW04]{AW04} S. Alaca, K. S. Williams, \emph{Introductory
  Algebraic Number Theory}, Cambridge University Press, 2004.
  (Standard modern account of pure cubic integral bases and unit
  inequalities.)

\bibitem[BG96]{BG96} Y. Bugeaud, K. Gy\H{o}ry, \emph{Bounds for the
  solutions of Thue--Mahler equations and norm form equations}, Acta
  Arith. 74 (1996). (Best general effective height bounds; polynomial
  in coefficient magnitude, hence exponential in coefficient
  bit-size.)

\bibitem[BH96]{BH96} Y. Bilu, G. Hanrot, \emph{Solving Thue equations
  of high degree}, J. Number Theory 60 (1996) 373--392.

\bibitem[BW87]{BW87} J. Buchmann, H. C. Williams, \emph{On principal
  ideal testing in algebraic number fields}, J. Symbolic Comput. 4
  (1987) 11--19. (Generator recovery
  $\mathfrak{a} = (\mu\,\alpha_k)$ along the cycle.)

\bibitem[BW88]{BW88} J. Buchmann, H. C. Williams, \emph{On the
  infrastructure of the principal ideal class of an algebraic number
  field of unit rank one}, Math. Comp. 50 (1988) 569--579. (The
  deterministic unconditional rank-one toolkit this paper runs on.)

\bibitem[Ded00]{Ded00} R. Dedekind, \emph{\"Uber die Anzahl der
  Idealklassen in reinen kubischen Zahlk\"orpern}, J. reine angew.
  Math. 121 (1900) 40--123. (The pure cubic index formula of
  Section~\ref{sec:index}.)

\bibitem[Hal07]{Hal07} S. Hallgren, \emph{Polynomial-time quantum
  algorithms for Pell's equation and the principal ideal problem},
  J. ACM 54 (2007). (The quantum route to the unit; cited in
  Remark~\ref{rem:beyond}.)

\bibitem[Lag79]{Lag79} J. C. Lagarias, \emph{Succinct certificates
  for the solvability of binary quadratic Diophantine equations},
  Proc. 20th IEEE FOCS (1979) 47--54; extended version
  arXiv:math/0611209. (\S 1.4 of the extended version quotes Smale's
  $2^{s^c}$ formulation.)

\bibitem[Len92]{Len92} H. W. Lenstra, Jr., \emph{Algorithms in
  algebraic number theory}, Bull. Amer. Math. Soc. 26 (1992)
  211--244. (Thm. 6.5: $hR \le |\Delta|^{1/2}\mathrm{polylog}$ and
  the ideal-count bound; Thm. 5.5 cited as hedged fallback only.)

\bibitem[Mat00]{Mat00} E. M. Matveev, \emph{An explicit lower bound
  for a homogeneous rational linear form in the logarithms of
  algebraic numbers. II}, Izv. Math. 64:6 (2000) 1217--1269.
  (Corollary 2.3, as restated in \cite[\S 2.4]{Sha19}.)

\bibitem[MST84]{MST84} M. Mignotte, T. N. Shorey, R. Tijdeman,
  \emph{The distance between terms of an algebraic recurrence
  sequence}, J. reine angew. Math. 349 (1984) 63--76.

\bibitem[PARI]{PARI} The PARI Group, \emph{PARI/GP version 2.17.4},
  Univ. Bordeaux. (\texttt{thueinit}/\texttt{thue} --- flag-1
  certification unconditional; \texttt{bnfinit},
  \texttt{bnfcertify}, \texttt{bnfisintnorm}.)

\bibitem[Sch08]{Sch08} R. Schoof, \emph{Computing Arakelov class
  groups}, in: Algorithmic Number Theory (J. P. Buhler,
  P. Stevenhagen, eds.), MSRI Publications 44, Cambridge University
  Press, 2008, 447--495. (General-rank infrastructure context for
  Sections~\ref{rem:general}--\ref{rem:totallyreal}.)

\bibitem[Sha19]{Sha19} M. Sha, \emph{Effective results on the Skolem
  Problem for linear recurrence sequences}, J. Number Theory 197
  (2019) 228--249; arXiv:1505.07147. (Theorems 1.1--1.2, Lemma 2.8,
  and the Matveev application --- the extracted ingredients of
  Lemma C.)

\bibitem[Sma96]{Sma96} N. P. Smart, \emph{How difficult is it to
  solve a Thue equation?}, in: Algorithmic Number Theory (ANTS-II),
  Lecture Notes in Computer Science 1122, Springer, 1996, 363--373.
  (The closest prior: per-method complexity estimate for the
  Tzanakis--de Weger pipeline.)

\bibitem[Sma98]{Sma98} S. Smale, \emph{Mathematical problems for the
  next century}, Math. Intelligencer 20:2 (1998) 7--15; reprinted in
  \emph{Mathematics: Frontiers and Perspectives}, AMS, 2000.

\bibitem[TdW89]{TdW89} N. Tzanakis, B. M. M. de Weger, \emph{On the
  practical solution of the Thue equation}, J. Number Theory 31
  (1989) 99--132.

\bibitem[Ver85]{Ver85} N. K. Vereshchagin, \emph{Occurrence of zero
  in a linear recursive sequence}, Math. Notes Acad. Sci. USSR 38:2
  (1985) 609--615.

\bibitem[Wil85]{Wil85} H. C. Williams, \emph{Continued fractions and
  number-theoretic computations}, Rocky Mountain J. Math. 15 (1985)
  621--655. (Venue sometimes misquoted as Pacific J. Math.; the
  correct venue is Rocky Mountain J. Math.)

\bibitem[Wil86]{Wil86} H. C. Williams, \emph{The spacing of the
  minima in certain cubic lattices}, Pacific J. Math. 124 (1986)
  483--496. (Spacing constants; scope: all complex cubic fields and
  arbitrary reduced lattices.)

\end{thebibliography}

\end{document}
